\documentclass[conference]{IEEEtran}
\usepackage[T1]{fontenc}
\usepackage{graphicx,booktabs,amsmath,amssymb,xcolor,url,multirow}
\usepackage[hidelinks]{hyperref}
\usepackage{cite}
\usepackage{xr-hyper}
\externaldocument[M-]{main}
% generated by scripts/paper_numbers.py -- do not edit
\providecommand{\TBD}{\textbf{\textcolor{red}{TBD}}}
\newcommand{\CN}{30}
\newcommand{\CVlmCorrect}{28}
\newcommand{\DN}{42}
\newcommand{\DCorrect}{42}
\newcommand{\DCorrectCI}{[91.6, 100.0]\%}
\newcommand{\DWrong}{0}
\newcommand{\DWrongCI}{[0.0, 8.4]\%}
\newcommand{\DLatMed}{3.05}
\newcommand{\DLatMax}{6.64}
\newcommand{\DRetries}{2}
\newcommand{\DWallMed}{102}
\newcommand{\EN}{2}
\newcommand{\ECorrect}{2}
\newcommand{\ECorrectCI}{[34.2, 100.0]\%}
\newcommand{\EWrong}{0}
\newcommand{\EWrongCI}{[0.0, 65.8]\%}
\newcommand{\ELatMed}{3.34}
\newcommand{\ELatMax}{3.76}
\newcommand{\ERetries}{0}
\newcommand{\EWallMed}{116}
\newcommand{\ESafeN}{1}
\newcommand{\ETaskSuccess}{1}
\newcommand{\ETaskSuccessCI}{[20.7, 100.0]\%}
\newcommand{\EBzeroWrong}{2}
\newcommand{\EBzeroWrongCI}{[34.2, 100.0]\%}
\newcommand{\EBzeroTask}{0}
\newcommand{\EBoneWrong}{0}
\newcommand{\EBoneTask}{0}
\newcommand{\IntN}{60}
\newcommand{\IntCorrect}{55}
\newcommand{\IntUnsafe}{0}
\newcommand{\IntTfidfCorrect}{31}
\newcommand{\IntTfidfUnsafe}{2}
\newcommand{\TNimages}{600}
\newcommand{\TNamb}{1}
\newcommand{\TNdeny}{408}
\newcommand{\TNallow}{191}
\newcommand{\TFA}{5}
\newcommand{\TFARate}{1.2\%}
\newcommand{\TFACI}{[0.5, 2.8]\%}
\newcommand{\TCA}{187}
\newcommand{\TCARate}{97.9\%}
\newcommand{\TCACI}{[94.7, 99.2]\%}
\newcommand{\THazRecall}{67/76}
\newcommand{\TConeAlarm}{0/25}
\newcommand{\TLatPass}{0.98}
\newcommand{\TAmbAllowed}{1}
\newcommand{\TSupplyN}{332}
\newcommand{\TSupplyFA}{0}
\newcommand{\TSupplyFACI}{[0.0, 1.1]\%}
\newcommand{\TFireCorrectN}{48}
\newcommand{\TFireCorrectFA}{5}
\newcommand{\TAblNoHazFA}{47}
\newcommand{\TvFA}{3}
\newcommand{\TvFARate}{0.7\%}
\newcommand{\TvFACI}{[0.2, 2.1]\%}
\newcommand{\TvCA}{187}
\newcommand{\TvCARate}{97.9\%}
\newcommand{\TvCACI}{[94.7, 99.2]\%}
\newcommand{\TvHazRecall}{69/76}
\newcommand{\TvConeAlarm}{0/25}
\newcommand{\TvNoFireAlarm}{0/523}
\newcommand{\TNoFireAlarm}{0/523}
\newcommand{\TimingN}{14}

\newcommand{\sys}{RescueMind}
\newcommand{\fig}[2][\linewidth]{\IfFileExists{#2}{\includegraphics[width=#1]{#2}}{\fbox{\parbox{0.9\linewidth}{\centering\TBD: \detokenize{#2}}}}}
\newcommand{\taus}{\tau_{s}}
\newcommand{\tauh}{\tau_{h}}
\renewcommand{\thesection}{S\arabic{section}}
\renewcommand{\thetable}{S\arabic{table}}
\renewcommand{\thefigure}{S\arabic{figure}}

\begin{document}
\title{Supplementary Material\\Verify, Look Again, Remember: Fail-Closed Vision-Language Checks, Active Look-Around and Safety-Monotone Mission Memory for Relief-Supply Delivery Robots}
\author{\IEEEauthorblockN{Anonymous Authors}\IEEEauthorblockA{Code, data and every trial log: \url{https://github.com/fe2131-art/rescue-mind}}}
\maketitle
This document holds material moved out of the main paper. Section, table and figure numbers with an S prefix refer to this
document; unprefixed references are to the main paper. All numbers are generated from the same result files as the main paper.

\section{Failure Analysis and Lessons}\label{sec:findings}
Running the full stack on a GPU revealed problems that \NTests{} offline unit tests and mocked ROS runs had not.
\begin{itemize}
\item \textbf{The auditor trusted the controller's sensor.} The wheel joint axis was expressed in the
child-link frame after a 90$^\circ$ roll, so the wheels spun about the vertical axis. Odometry reported 1.5\,m forward motion while the
body slid sideways by 0.75\,m. An audit based on \texttt{/odom} would have passed the run.
Audits now use Gazebo's ground-truth odometry. Changing the wheel collision from a cylinder to a sphere reduced
dead-reckoning drift from 0.16 to 0.011\,m.
\item \textbf{Free-form JSON from the first VLM was unusable} on our renders. With Qwen2.5-VL-3B, 21 of 48 outputs could not be parsed and the model hallucinated water on empty shelves.
A newer model parses reliably, but it still approves fewer correct deliveries than scoring does (Table~\ref{M-tab:t1models}).
This motivated closed-set scoring. T1 measures the free-JSON method again (Table~\ref{M-tab:t1models}).
\item \textbf{Middleware effects.} Best-effort QoS silently dropped 92 of 102 0.9\,MB camera frames under load. The GPU lidar
detected the robot's own camera mast. Slow DDS discovery on a 4-core host failed 4 of 12 starts in an aborted campaign.
An interrupted supervisor left an orphaned stack that made every trial of the next campaign stop with two
\texttt{/cmd\_vel} publishers. All of these are recorded with their fixes; the aborted campaigns are published.
\item \textbf{Stale-request latch.} Two pre-decision retries in campaign D happened while the robot was stationary: the mission node's
request stream stalled for $\approx$0.3\,s on the saturated CPU. We instrumented the mission loop, the guard
and garbage collection and rejected two hypotheses (guard starvation, GC pauses). Raising control-loop priority did not
remove the stalls. We therefore changed the guard: a stale zero request now holds zero, while a stale motion request still latches.
With a 0.6\,s stall injected before departure, the old guard failed 2/2 runs and the new one completed 2/2.
Raising priority with a file capability then had a side effect: it put the exec into the dynamic loader's secure mode,
which stripped \texttt{LD\_LIBRARY\_PATH} and broke the ROS bag recorder. A regression run caught this before any campaign.
\item \textbf{A pilot exposed a threshold that was not conservative enough.} Before campaign E we ran nine trials in
separate pilot worlds. In one, the robot found the wrong item, searched the other station, and delivered from it although a fire
burned behind it. The fire was almost entirely hidden by the first-aid kit, and $p_h=0.32$. On development data no fire-free
image exceeded $p_h=0.0007$, so we lowered $\tauh$ to 0.05 (v3.5). The same pilot worlds then gave 9/9 expected outcomes.
T1 was not involved: both thresholds were registered for T1 before any T1 image was scored.
\item \textbf{A generator bug caught by looking at the data.} In the first T1 render, distractors meant to stand next to a supply were placed through it. The
emptiness check ran after the supplies were moved. The render was discarded and the fix recorded as a
pre-registration amendment, before any model scored a T1 image. A second amendment replaced the field-of-view fire
label with the rendered-mask label, after we saw fires hidden behind supplies that the geometric label had called visible.
\end{itemize}

\section{Wall-clock timing}\label{sec:timing}
\textbf{Where the time goes.} Table~\ref{tab:timing} breaks down the \TimingN{} delivered missions that needed no search.
The check itself takes about \TimCheck\,s: aiming, waiting until stationary, frame checks, and the scored VLM queries
(\TimVlm\,s). Driving between station and shelter takes about \TimDrive\,s (sum of phase medians). The check is therefore short
compared with driving. The situation report adds about 4\,s but could run in parallel with the drive to the shelter.
\begin{table}[t]\centering\caption{Wall-clock phases of delivered missions without search (campaigns D and E; v3.4 and v3.5, no look-around).}\label{tab:timing}
\resizebox{0.92\linewidth}{!}{% generated by scripts/paper_numbers.py
\begin{tabular}{lrrr}
\toprule
Phase & Median [s] & Min & Max \\
\midrule
drive to station & 14.31 & 12.19 & 16.62 \\
aim camera & 0.50 & 0.50 & 0.52 \\
wait until stationary & 0.67 & 0.66 & 0.71 \\
acquire and check frames & 0.90 & 0.75 & 0.98 \\
VLM decision (scored queries) & 3.44 & 2.95 & 6.64 \\
situation report (not used to decide) & 5.14 & 3.99 & 8.54 \\
drive to shelter & 24.83 & 15.02 & 25.32 \\
\bottomrule
\end{tabular}
}\end{table}

\section{Additional T1 analyses}\label{sec:t1extra}
\subsection{Decision-rule ablations}
\begin{table}[t]\centering\caption{Decision-rule ablations on T1: false approvals / false denials.
$^{*}$: exact McNemar $p<0.05$ versus the full rule of the same model.}
\label{tab:t1ablation}\resizebox{\linewidth}{!}{% generated by scripts/paper_numbers.py
\begin{tabular}{lr}
\toprule
Decision rule & Qwen3.5-9B \\
\midrule
Full rule & 5 / 4 \\
A1: question 1 only & 5 / 4 \\
A2: question 2 only & 5 / 1 \\
A3: no hazard question & 47 / 4$^{*}$ \\
A4: no threshold ($\tau_s{=}0$) & 5 / 4 \\
v3.5: $\tau_h{=}0.05$ (I-10) & 3 / 4 \\
\bottomrule
\end{tabular}
}\end{table}
Table~\ref{tab:t1ablation} isolates the parts of the rule on T1 (Table~\ref{M-tab:t1models}, Eq.~\ref{M-eq:rule}). For the 9B model neither the
permuted second question nor the confidence threshold changed the number of false approvals. With Qwen2.5-VL-3B, requiring the two
multiple-choice questions to agree reduced false approvals, which is the option-letter bias of \cite{zheng2024mcselectors}.
\subsection{Risk--coverage}
\begin{figure}[t]\centering\fig{fig_risk_coverage.pdf}
\caption{T1 risk--coverage: false-approval rate versus correct-approval rate as $\taus$ varies ($\tauh=0.05$, v3.5).
Markers show the frozen $\taus=0.5$. The curves are flat over most of the range: the remaining errors are missed fires,
which $\taus$ cannot fix.}\label{fig:riskcov}\end{figure}
Fig.~\ref{fig:riskcov} varies $\taus$ only for description, because $\taus$ was not tuned on T1.

\section{Campaign E details}\label{sec:camp-e}
Campaign E used a new randomised world in every trial (main paper, Sec.~\ref{M-sec:results}); Table~\ref{tab:e} gives the outcomes by stratum and
Fig.~\ref{fig:etraj} the ground-truth trajectories.
\begin{table}[t]\centering\caption{Campaign E (randomised worlds) by stratum.}\label{tab:e}
\resizebox{\linewidth}{!}{% generated by scripts/paper_numbers.py
\begin{tabular}{llrrl}
\toprule
Expected & Condition & $n$ & Correct & Observed \\
\midrule
deliver after search & plain & 1 & 1 & allow after search 1 \\
deny & crate & 1 & 1 & deny 1 \\
\bottomrule
\end{tabular}
}\end{table}
\begin{figure}[t]\centering\fig{fig_e_trajectories.pdf}
\caption{Campaign E: ground-truth trajectories of all final trials, coloured by audited outcome.}\label{fig:etraj}\end{figure}

\section{Model selection on development data}\label{app:models}
\begin{table}[h]\centering\caption{Candidate models with the final prompts (\texttt{v35}: checkpoint questions of v3.5; \texttt{i35}: instruction parser) on development data only.
Images: correct decisions with false approvals (FA) and false denials (FD). Explanations: one-sentence reports judged
correct against the generator labels. Intent: development utterances, with unsafe acceptances. VRAM in MiB.}
\label{tab:devmodels}\resizebox{\linewidth}{!}{% generated by scripts/paper_numbers.py
\begin{tabular}{llrrrr}
\toprule
Model & Released & Images (176) & Explanations (128) & Intent dev (40) & VRAM \\
\midrule
Qwen2.5-VL-3B & 2025-01 & 176/176 (FA 0/FD 0) & 84/128 & 31/40, unsafe 0 & -- \\
Qwen3-VL-8B & 2025-10 & 176/176 (FA 0/FD 0) & 120/128 & 38/40, unsafe 2 & 17212 \\
Qwen3.5-4B & 2026-02 & 175/176 (FA 1/FD 0) & 125/128 & 35/40, unsafe 2 & 9102 \\
Qwen3.5-9B & 2026-02 & 176/176 (FA 0/FD 0) & 125/128 & 35/40, unsafe 0 & 18320 \\
\bottomrule
\end{tabular}
}\end{table}
Models were ranked lexicographically. The first key was zero false approvals and zero unsafe intent acceptances, and the
second was explanation and intent accuracy. Qwen3-VL-8B parsed more development utterances but accepted two unsafe
ones. Qwen3.5-4B produced a false approval. Qwen3.5-9B was therefore chosen and frozen before the T1 and E registrations.

\section{Registration audit}\label{app:prereg}
\begin{table}[h]\centering\caption{When each study was registered and when its first results entered the repository. ``Order''
checks that the registration commit is an ancestor of the first results commit (\texttt{git merge-base --is-ancestor}).
Registrations are also mirrored in the public repository history.}\label{tab:prereg}
\resizebox{\linewidth}{!}{% generated by scripts/paper_numbers.py
\begin{tabular}{llllc}
\toprule
Study & Registered (commit, date) & First results (commit, date) & Amendments & Order \\
\midrule
C (v3.3) & \texttt{5c9a951 2026-10-16} & \texttt{5c9a951 2026-10-16} & 0 & same commit \\
D (v3.4) & \texttt{ac6d815 2026-10-16} & \texttt{fb4e9b3 2026-10-16} & 0 & $\checkmark$ \\
T1 & \texttt{5f835af 2026-10-16} & \texttt{c3d1ffc 2026-10-16} & 3 & $\checkmark$ \\
E & \texttt{6201aa0 2026-10-16} & \texttt{17e0ea6 2026-10-16} & 0 & $\checkmark$ \\
T2 & \texttt{f347239 2026-10-17} & \texttt{7843e9c 2026-10-17} & 1 & $\checkmark$ \\
F & \texttt{564a724 2026-10-17} & \texttt{5bbf11a 2026-10-17} & 0 & $\checkmark$ \\
G & \texttt{edb26a1 2026-10-17} & \texttt{f6160fb 2026-10-17} & 1 & $\checkmark$ \\
R & \texttt{fb63bee 2026-10-17} & \texttt{pending} & 0 & pending \\
H & \texttt{a5cc651 2026-10-17} & \texttt{pending} & 1 & pending \\
J & \texttt{b72b877 2026-10-17} & \texttt{pending} & 0 & pending \\
Z & \texttt{a53ee7e 2026-10-17} & \texttt{pending} & 0 & pending \\
\bottomrule
\end{tabular}
}\end{table}
For every study with results, the registration commit precedes the results commit (\PreregOrderOK). Campaign C's
registration and ledger entered the repository together, so their order cannot be verified from the history.


\bibliographystyle{IEEEtran}
\bibliography{references}
\end{document}
